\worksheettitle
  {Ноль и разрядный состав}
  {\modulemark{M02} Версия учителя}

\begin{teacherbox}[Паспорт модуля]
  \textbf{Результат:} ученик объясняет роль нуля в позиционной записи,
  раскладывает число на разрядные слагаемые и восстанавливает число по сумме.

  \textbf{Предварительно:} ученик определяет разряд и значение цифры (M01).

  Модуль не ограничен временем. Обязательный маршрут заканчивается контрольной
  точкой; объём практики выбирается по устойчивости способа.
\end{teacherbox}

\begin{checkpointbox}[Вход: ответы]
  Значения цифр \(8\): \(80\,000\) и \(800\). Нули стоят в разрядах единиц
  тысяч и десятков. Если ученик называет только цифры или считает разряды
  слева, сначала используйте M01-R.
\end{checkpointbox}

\section{Открытие роли нуля}

В числе \(40\,508\) нет единиц тысяч и десятков. В таблице следует записать
нули и значения \(0\). В числе \(458\) цифра \(4\) имеет значение \(400\), а
цифра \(5\) \textemdash{} \(50\). В числе \(40\,508\) их значения равны \(40\,000\) и
\(500\).

\begin{teacherbox}[Ключевой смысл]
  Не ограничивайтесь фразой «ноль нельзя убрать». Ученик должен связать ноль с
  отсутствием единиц определённого разряда и изменением значений остальных
  цифр при сдвиге записи.
\end{teacherbox}

При удалении нуля в разряде единиц тысяч из \(60\,205\) получится \(6\,205\).
При удалении нуля в разряде десятков получится \(6\,025\). В обоих случаях
значение цифры \(6\) уменьшается с \(60\,000\) до \(6\,000\), то есть на
\(54\,000\). Во втором случае значение цифры \(2\) уменьшается с \(200\) до
\(20\), то есть на \(180\). Следовательно, больше изменение значения цифры
\(6\). Ученик сравнивает величины изменений \(54\,000\) и \(180\), а не
пользуется неопределённой формулировкой «изменилось сильнее».

\section{Разложение: ответы}

Пропуски:
\begin{enumerate}
  \item \(3\,000\) и \(4\);
  \item \(500\);
  \item \(5\,000\).
\end{enumerate}

Самостоятельное разложение:
\begin{align*}
  63\,204&=60\,000+3\,000+200+4,\\
  8\,090&=8\,000+90,\\
  50\,006&=50\,000+6,\\
  304\,080&=300\,000+4\,000+80,\\
  600\,007&=600\,000+7,\\
  701\,040&=700\,000+1\,000+40.
\end{align*}

В числе \(304\,080\) в разрядах десятков тысяч, сотен и единиц стоят нули,
поэтому соответствующие нулевые слагаемые обычно не записывают.

Проверочная пара:
\[
  70\,305=70\,000+300+5,
  \qquad
  70\,000+300+5=70\,305.
\]

\section{Сборка: ответы}

В таблице для \(80\,000+700+5\) отсутствуют единицы тысяч и десятков: в этих
разрядах записываются нули. Получается \(80\,705\).

\begin{enumerate}
  \item \(53\,204\);
  \item \(80\,705\);
  \item \(340\,008\);
  \item \(600\,021\);
  \item \(904\,060\);
  \item \(700\,005\).
\end{enumerate}

\begin{teacherbox}[Два направления \textemdash{} две ошибки]
  При разложении ученик может неверно определить значение цифры. При сборке он
  чаще пропускает отсутствующий разряд. Диагностируйте направление действия,
  а не только конечный ответ.
\end{teacherbox}

\section{Ошибки: ответы и причины}

\begin{enumerate}
  \item \(80\,000+700+5=80\,705\). Потеряны разряды единиц тысяч и десятков.
  \item \(8\,090=8\,000+90\). Цифра \(9\) стоит в разряде десятков, а не сотен.
  \item \(50\,204=50\,000+200+4\). Цифре \(2\) ошибочно приписано значение
    \(2\,000\); кроме того, значение \(200\) записано второй раз.
  \item \(300\,000+4\,000+8=304\,008\). При сборке потеряны нулевые позиции.
\end{enumerate}

\section{Конструирование}

Возможные ответы:
\begin{itemize}
  \item пятизначное число без единиц тысяч и десятков: \(50\,304\);
  \item шестизначное число с тремя ненулевыми слагаемыми: \(405\,070\),
    поскольку \(405\,070=400\,000+5\,000+70\).
\end{itemize}
Принимайте любые числа, удовлетворяющие условиям, если ученик верно проверил
их разложением.

При вставке нулей в \(428\) возможны разные ответы. Например, \(4\,028\) и
\(40\,028\). Оценивайте сохранение порядка цифр \(4,2,8\), требуемую длину
числа и правильное определение значений цифр.

\section{Контрольная точка}

\begin{align*}
  90\,307&=90\,000+300+7,\\
  500\,040&=500\,000+40,\\
  30\,000+400+8&=30\,408,\\
  600\,000+2\,000+70&=602\,070.
\end{align*}
В первом числе в разрядах единиц тысяч и десятков стоят нули, поэтому
соответствующие нулевые слагаемые не записаны.

\begin{resultbox}[Решение о маршруте]
  \begin{itemize}
    \item \textbf{Переход к M03:} четыре преобразования выполнены верно,
      ученик объясняет отсутствие нулевых слагаемых через разряды.
    \item \textbf{M02-R:} потерян разряд при сборке, цифре приписано неверное
      значение или объяснение роли нуля отсутствует.
    \item \textbf{Дополнительная практика или M02-H:} способ понятен, но в
      длинных числах появляются единичные ошибки.
    \item \textbf{M02\(\star\):} базовый результат устойчив; можно переходить к
      смешанным разрядным единицам.
  \end{itemize}
\end{resultbox}

\begin{teacherbox}[После коррекции]
  Дайте новую проверку: разложить \(80\,406\) и собрать
  \(500\,000+3\,000+20+7\). Ожидаются
  \(80\,000+400+6\) и \(503\,027\).
\end{teacherbox}

\newpage
\section{Ответы к M02-R}

\begin{teacherbox}[Пошаговое восстановление способа]
  В таблице для \(60\,000+3\,000+40+2\) в разряде сотен записываются значение
  \(0\) и цифра \(0\); число равно \(63\,042\). Ноль сохраняет место разряда
  сотен и позиции остальных цифр.

  Для \(80\,507\): цифры \(8,0,5,0,7\), значения
  \(80\,000,0,500,0,7\), разложение \(80\,000+500+7\).
  Короткая проба: \(70\,604\); \(90\,506=90\,000+500+6\).
\end{teacherbox}

\begin{teacherbox}[Ошибки и повторная проверка]
  В числе \(70\,506\) цифра \(5\) стоит в разряде сотен и имеет значение
  \(500\); верно \(70\,506=70\,000+500+6\).

  В записи \(80\,000+700+5=875\) потеряны разряды единиц тысяч и десятков;
  верное число \(80\,705\).

  Исправления: \(60\,090=60\,000+90\) и
  \(400\,000+3\,000+8=403\,008\).

  Повторная проверка: \(40\,509\) и
  \(70\,008=70\,000+8\). В первом числе нули сохраняют разряды единиц тысяч
  и десятков; во втором нулевые слагаемые тысяч, сотен и десятков не
  записываются.
\end{teacherbox}

\newpage
\section{Ответы к M02\(\star\)}

\begin{teacherbox}[Переводим разрядные единицы]
  Разобранный пример:
  \[
    6\cdot10\,000=60\,000,\qquad
    34\cdot100=3\,400,\qquad
    60\,000+3\,400+8=63\,408.
  \]

  Самостоятельная практика:
  \[
    42\,705;\qquad 325\,140;\qquad 86\,032.
  \]
\end{teacherbox}

\begin{teacherbox}[Сравнение и обратные задачи]
  Сравнения: \(63\,400=63\,400\); \(502\,800\ne502\,080\);
  \(70\,160=70\,160\).

  Количество полных сотен в числе \(4\,700\) равно \(47\), количество полных
  тысяч в \(62\,000\) равно \(62\), количество полных десятков в \(3\,540\)
  равно \(354\).
  Представление \(28\,600\): \(2\) десятка тысяч, \(86\) сотен и \(0\)
  единиц, то есть \(20\,000+8\,600=28\,600\).
\end{teacherbox}

\begin{teacherbox}[Проверка чужого решения и обобщение]
  \(4\) десятка тысяч и \(27\) сотен дают
  \(40\,000+2\,700=42\,700\), а не \(427\): записи количеств нельзя просто
  приписывать друг к другу.

  Собственный пример принимается, если все названные части переведены в
  числовые значения, сумма вычислена верно и в полученном числе действительно
  есть нулевой разряд. В обобщении ученик объясняет, что одинаковые числа
  разрядных единиц могут относиться к разным разрядам, поэтому сначала нужно
  найти значение каждой части.
\end{teacherbox}
